\subsection{Problem Statement}
The main problem that this thesis aims to solve comes from a consideration while thinking how AI can be implemented in a blue team scenario.
Even if this topic could be an argument of discussion by itself, many perplexities arise when we think how data is exchanged considering AI agents actively inside the attack surface.

In particular, the problem is that AI agents are not able to distinguish between sensitive and non-sensitive data, and they are not able to understand the context of the data they are processing. This means that if an AI agent is given access to sensitive data, it could potentially leak that data to an attacker, either intentionally or unintentionally.
Of course, it is imperative to also consider the intrinsic vulnerability of AI agents to prompt injection attacks, which could be just another mean to the exfiltration of sensitive data as upmentioned.

One method to mitigate this problem would be to utilize cryptography, but that alone comes with its own set of challenges. More in detail, both an AI agent and a user cannot benefit from data that is encrypted with a symmetric key: it is easy to see how both of them would require a decryption of said  data, otherwhise they would not be able even to understand it, leave alone to use it.

One possible alternative could be Homomorphic Encryption and, surely, it is a very promising topic, but as for today its computational limitations make it not suitable for a real-world scenario, especially if we consider the amount of data that an AI agent should process.

Another approach to achieve more security in a scenario that involves a LLM could be Trusted Execution Environments, but here lies the need (thus the problem) to ensure that the AI agent is running in a trusted environment, which is not always possible, especially if we consider that the AI agent could be running on a remote server or in a cloud environment.

Hereby we'll expand on the limitations of both Homomorphic Encryption and Trusted Execution Environments, to better understand and define the problem we are trying to solve.
\subsubsection{Limitations of Homomorphic Encryption}
As upmentioned, Homomorphic Encryption is a very promising topic, given the ability it enables to perform homomorphic operations on ecnrypted data. The problems arise when we approach statistical computations, which is the basis of Large Language Models, and, in general, of any AI agent.

Precisely on this regard, \citeauthor{cheon2017homomorphic} eloquently stated \cite{cheon2017homomorphic}:
\begin{quote}
    Most of real-world data contain some errors from their true values. [...] an approximate value may substitute the original data and a small rounding
error does not have too much effect on computation result [...] an approximate value may substitute the original data and a small rounding
error does not have too much effect on computation result [...]
    Unfortunately this rounding operation has been considered difficult to perform on HE
since it is not simply represented as a small-degree polynomial.
\end{quote}
They proposed a scheme that allows to perform approximate computations on encrypted data, not without open problems, like the need to convert said scheme in a fully homomorphic scheme using bootstrapping. This problem was later addressed by \citeauthor{de2025encryptedllm} who proposed a GPU-based workflow by extending the capabilities of OpenFHE \cite{al2022openfhe}.

Said implementations is undoubtly a great step forward towards practical encrypted LLM, but it has still non trivial limitations:
\begin{quote}
    Although promising given the massive overheads involved in both LLMs
    and FHE, this is still far from practical for real-world usage,
    even for applications that are not latency sensitive such
    as text summarization or content generation (in contrast to
    chatbots or Q/A tasks, which are more demanding in terms
    of responsiveness).
\end{quote}
The main problem is that, even with a GPU-accelerated implementation, the bootstrap overhead (also in combination with the LLM inference overhead) is still too high to be considered practical in a every day scenario.
To this we have to take in account also that by approximating the functions needed to have LLM inference we are introducing a source of error that translates into a loss of accuracy in the final result.
Both of this factors alone make Homomorphic Encryption, as per today, not viable to be used in application such as chatbot or coding agents and it is clear why Attribute-Based Encryption could fill the gap in this point and time by providing security precisely where the FHE, at the moment, cannot provide.
\subsubsection{Limitations of Trusted Execution Environments}
Trusted Execution Environmets, as the name suggests, are trusted environment where it is possible to execute critical operation and, in the context of prompt injection and LLM, potentially separate a network or an entire device from the execution of malicious actions by the agent that, in fact, would only be relegated in a sandbox.

\citeauthor{geppert2022trusted} define TEEs as following:
\begin{quote}
    TEEs [...] are defined as tamper-resistant processing environments that run on a separation kernel.
    Remote verification makes the technology suitable for cloud
    environments. It is a mechanism proving to
    an organization that uses the services of a cloud provider that the
    specified software runs on the predefined hardware even if the
    organization does not have direct access to this hardware.
\end{quote}

In the cloud computing scenario this sort of infrastruction is already offered by providers like Google, Amazon, and Microsoft \cite{geppert2022trusted}, but this would not apply to some organization that don't want to rely on external provider. This is extremely crucial even when talkin of LLMs, since if an organization opted to rely on third party agents (i.e. Claude Code, Codex, etc...) it would rely to this third party also all the security assumtion that arise.

To better explain this concept, by considering the use of a TEE provided by some provider, an organization would not have any control on \textit{how} its own data are manipulated. It is obvious that in most cases this is not an acceptable risk and many would resort to a "home-made" solution, like training its own LLM; this issue, although, could be mitigated if we consider a situation where only pre-approved data can be "sent out", thus the application of Attribute-Based Encryption and the not suitability of TEE in every scenario.

\subsection{Actors and Goals}
In this section we define the actors involved in the system and their goals, to better understand the problem we are trying to solve.

The actors, with their own goals, involved in the system are:
\begin{itemize}
    \item \textbf{User:} the user is the person that wants to use an AI agent to process a document. The user has a set of attributes that define their access rights to the document.
    \item \textbf{AI Agent:} the AI agent is the entity that retrieves and processes the document. The AI agent has a set of attributes that define its access rights to the document.
    \item \textbf{Document:} the document is the entity that contains sensitive data. The document is encrypted with a CP-ABE scheme and has an access policy that defines who can access it.
    \item \textbf{Company:} the company is the entity that owns the document and defines the access policy for it alongside the attributes of the users and the AI agents. The company is responsible for ensuring that the access policy is enforced and that the sensitive data is protected.
\end{itemize}
% TODO: add the use case UC1 to better explain the actors and goals. (Maybe better in System Design)

\subsection{Assumptions}
We assume that the documents subject to protection are already produced and encrypted via Attribute-Baed Encryption before being made available to the system described in this work. Authoring a new document, defining its access policy, ensuring the freshness of employees' attributes, generating the corresponding keys, encrypting the document, are all separate concerns and fall out of the scope of this thesis, which focuses on enforcing, at read time, a mechanism that allows an AI agent to process the document following a safe workflow, minimizing the risk of sensitive data exfiltration.

We also assume that our garantees end at the moment decrypted content is legitimately handed to the requesting AI agent, in accordance with the CP-ABE policy and the combined attributes of the user and the agent.
What the agent does with that content afterwards is not something an access-control layer operating at application level can controll, and is therefore out of scope. There is an exception worth to mention, that is locally materialized copies which are time-bounded and removed automatically at the end of a session, to limit possible unbounded exposure.

\subsection{Adversary Model}

\subsection{Security Goals}
